En este informe pondremos a prueba dos problemas clásicos de la ciencia de la computación: el ordenamiento de arreglos unidimensionales y la multiplicación de matrices cuadradas. Para la primera parte, implementamos los algoritmos Quick Sort, Merge Sort y Patience Sort, sumando también a \texttt{std::sort} (que ya venía listo en la base de la tarea). Por el lado de las matrices, enfrentamos al método tradicional (Naive) contra el algoritmo de Strassen. 

El objetivo principal de todo este análisis es llevar la teoría a la vida real: queremos contrastar el comportamiento empírico de estas implementaciones, midiendo cuánto tiempo se demoran y cuánta memoria consumen, con su complejidad teórica asintótica. De esta forma, buscaremos comprobar si los modelos matemáticos realmente se sostienen cuando los llevamos a la práctica.